\chapter{Memory Map and Bus Access}
\label{ch:memory-map}

\section{The Device Segment}

The sound device occupies one segment of the CPU's address space, 65,536 words long. The segment is chosen by the
system; the segment that the SIRCUS system uses is not yet specified. This manual gives every address as an offset from
the base of the segment.

\begin{table}[H]
	\centering
	\begin{tabular}{llp{6.5cm}}
		\toprule
		\textbf{Offsets}                 & \textbf{Region}     & \textbf{Notes}                                                                    \\
		\midrule
		\texttt{0x0000}--\texttt{0x000F} & Global registers    & Offsets \texttt{0x0000}--\texttt{0x000B} are defined (Chapter~\ref{ch:registers}) \\
		\texttt{0x0010}--\texttt{0x001F} & Voice 0 registers   & Eight registers defined (Section~\ref{sec:voice-registers})                       \\
		\texttt{0x0020}--\texttt{0x002F} & Voice 1 registers   & As voice 0                                                                        \\
		\texttt{0x0030}--\texttt{0x003F} & Voice 2 registers   & As voice 0                                                                        \\
		\texttt{0x0040}--\texttt{0x004F} & Voice 3 registers   & As voice 0                                                                        \\
		\texttt{0x0050}--\texttt{0x005F} & Voice 4 registers   & As voice 0                                                                        \\
		\texttt{0x0060}--\texttt{0x006F} & Voice 5 registers   & As voice 0                                                                        \\
		\texttt{0x0070}--\texttt{0x7FFF} & Reserved            & Reads zero, writes ignored                                                        \\
		\texttt{0x8000}--\texttt{0xFFFF} & Audio memory window & 32,768 words (Section~\ref{sec:audio-memory})                                     \\
		\bottomrule
	\end{tabular}
	\caption{Sound Device Memory Map}
	\label{tab:memory-map}
\end{table}

The \textbf{registers} control the device. The global registers are at offsets \texttt{0x0000}--\texttt{0x000F}. Each
of the six voices has a block of 16 registers, at \texttt{0x0010} plus 16 times the voice number. They are described in
Chapter~\ref{ch:registers}.

\section{The Audio Memory}
\label{sec:audio-memory}

The audio memory holds everything the voices play and everything that controls them: sample data, the instrument table,
the pitch tables and the sequences. It is 32,768 words, which is 64 KB. The device does not interpret the memory in any
fixed layout. Registers and instrument fields hold the addresses of the things they use, and any range of the memory
may hold any of them.

The CPU reaches the memory through the window at offsets \texttt{0x8000}--\texttt{0xFFFF}. The offset \texttt{0x8000}
plus $n$ selects word $n$ of the memory. The memory is not on the CPU's bus, and the CPU has no other way to reach it
(Section~\ref{sec:memory-interface}).

\subsection{Word and Byte Addresses}
\label{sec:byte-offsets}

Two kinds of address name the same memory, and a field says which it holds:

\begin{itemize}
	\item A \textbf{window address} is the CPU's word address of a word: \texttt{0x8000} for the first word, \texttt{0xFFFF} for
	      the last. \texttt{INSTR\_BASE} and \texttt{SEQ\_PC} hold window addresses.
	\item A \textbf{byte offset} is a count of bytes from the start of the audio memory, \texttt{0x0000}--\texttt{0xFFFF}. The
	      sample start, the pitch table address and the loop start of an instrument are byte offsets, as are the sample position
	      of a voice and its length. The byte at offset $b$ is in the word at window address $\texttt{0x8000} + \lfloor b/2
		      \rfloor$.
\end{itemize}

The sequencer's \texttt{GOTO} displacement is a count of instruction words (Chapter~\ref{ch:sequencer}).

\subsection{Byte Order}
\label{sec:byte-order}

A word holds two bytes, and the byte at the even offset is the high byte. For the word \texttt{0x1234}, the byte at an
even offset is \texttt{0x12} (bits 15--8) and the byte at the next odd offset is \texttt{0x34} (bits 7--0). Sample
bytes are signed two's-complement values. Multi-byte values in the tables of Chapter~\ref{ch:formats} are stored with
the high word first.

\section{Access Rules}
\label{sec:access-rules}

A read of a register, a reserved offset or a word of the window returns the stored value, or zero for a reserved
offset. A write stores the value, or is ignored at a reserved offset. No register has a side effect on a read. The
registers that act on a write are described with their registers (Chapter~\ref{ch:registers}).

\begin{itemize}
	\item \textbf{Registers} are accessed without delay: the device completes the cycle at once. A register access never
	      waits for the audio memory.
	\item \textbf{The window} is accessed in the device's memory schedule (Chapter~\ref{ch:timing}). A window access
	      is held until one of the slots reserved for the CPU comes round, and takes between two and five cycles to
	      complete, whatever the voices and sequencers are doing.
	\item \textbf{Playing and the window.} The CPU may read and write the audio memory while the voices play. The sound
	      processor reads in the first slot of each pattern of four slots and the CPU's window uses the other three, so a
	      read by the sound processor sees every write by the CPU that completed before its slot began. A program that
	      rewrites a word a voice is reading must not depend on whether the old or the new value is read for the one sample
	      in which the two coincide.
\end{itemize}

At reset the contents of the audio memory are undefined. A program must write every word it uses before it starts a
voice or a sequence that uses it (Section~\ref{sec:startup}).

\section{Planning an Audio Memory Layout}
\label{sec:layout}

Nothing fixes where the tables live, so a program chooses. One arrangement puts the instrument table at the start of
the memory, then the pitch tables, then the sequences, and gives the rest to sample data:

\begin{table}[H]
	\centering
	\begin{tabular}{lll}
		\toprule
		\textbf{Contents}                       & \textbf{Window addresses}  & \textbf{Size}                     \\
		\midrule
		Instrument table (\texttt{INSTR\_BASE}) & \texttt{0x8000} and up     & 8 words (16 bytes) per instrument \\
		Pitch tables                            & after the instrument table & 24 words (48 bytes) per table     \\
		Sequences                               & after the pitch tables     & One word per instruction          \\
		Sample data                             & the rest                   & One byte per sample               \\
		\bottomrule
	\end{tabular}
	\caption{An Example Audio Memory Layout}
	\label{tab:layout}
\end{table}

A sample takes one byte per sample point, so the whole memory holds at most about two seconds of sound at the device's
output rate if it is played at its original pitch with one sample point per output sample. Looping a short waveform
repeats its bytes, so a looped voice uses very little memory (Chapter~\ref{ch:playback}).
